% =========================================================
% 04_vector_spaces_basis.tex
% Vectors and Linear Algebra
% =========================================================

% =========================================================
% SECTION DIVIDER
% =========================================================

\section{Vector Spaces, Basis and Dimension}

\begin{frame}{Basis of a Vector Space}

A set of vectors

\[
\mathcal{B}
=
\{\vect{v}_1,\vect{v}_2,\ldots,\vect{v}_n\}
\]

is a \textbf{basis} of a vector space if:

\begin{enumerate}
    \item the vectors span the space, and
    \item the vectors are linearly independent.
\end{enumerate}

Therefore,

\[
\boxed{
\text{Basis}
=
\text{Spanning set}
+
\text{Linear independence}
}
\]

\vspace{0.4cm}

A basis provides a set of fundamental directions from which
every vector in the space can be constructed.

\end{frame}


% =========================================================
% STANDARD BASIS
% =========================================================

\begin{frame}{Standard Basis}

The standard basis of $\mathbb{R}^2$ is

\[
\boxed{
\mathcal{B}
=
\left\{
\begin{bmatrix}
1\\
0
\end{bmatrix},
\begin{bmatrix}
0\\
1
\end{bmatrix}
\right\}.
}
\]

Every vector

\[
\vect{x}
=
\begin{bmatrix}
x_1\\
x_2
\end{bmatrix}
\]

can be written as

\[
\vect{x}
=
x_1
\begin{bmatrix}
1\\
0
\end{bmatrix}
+
x_2
\begin{bmatrix}
0\\
1
\end{bmatrix}.
\]

\vspace{0.3cm}

Thus, the basis vectors provide the coordinate directions of
$\mathbb{R}^2$.

\end{frame}


% =========================================================
% BASIS IN R3
% =========================================================

\begin{frame}{Standard Basis of $\mathbb{R}^3$}

The standard basis of $\mathbb{R}^3$ is

\[
\mathcal{B}
=
\left\{
\begin{bmatrix}
1\\0\\0
\end{bmatrix},
\begin{bmatrix}
0\\1\\0
\end{bmatrix},
\begin{bmatrix}
0\\0\\1
\end{bmatrix}
\right\}.
\]

Any vector

\[
\vect{x}
=
\begin{bmatrix}
x_1\\x_2\\x_3
\end{bmatrix}
\]

can be expressed as

\[
\boxed{
\vect{x}
=
x_1\vect{e}_1
+
x_2\vect{e}_2
+
x_3\vect{e}_3.
}
\]

The three basis vectors provide three independent directions.

\end{frame}


% =========================================================
% COORDINATES RELATIVE TO A BASIS
% =========================================================

\begin{frame}{Coordinates Relative to a Basis}

Let

\[
\mathcal{B}
=
\{\vect{v}_1,\vect{v}_2\}
\]

be a basis of $\mathbb{R}^2$.

A vector $\vect{x}$ can be expressed as

\[
\vect{x}
=
c_1\vect{v}_1+c_2\vect{v}_2.
\]

The coefficients form the \textbf{coordinate vector} of $\vect{x}$
relative to $\mathcal{B}$:

\[
\boxed{
[\vect{x}]_{\mathcal{B}}
=
\begin{bmatrix}
c_1\\
c_2
\end{bmatrix}.
}
\]

\vspace{0.3cm}

The same geometric vector can therefore have different coordinates
when different bases are used.

\end{frame}


% =========================================================
% DIMENSION
% =========================================================

\begin{frame}{Dimension}

The \textbf{dimension} of a vector space is the number of vectors
in any basis of that space.

For example,

\[
\boxed{
\dim(\mathbb{R}^2)=2
}
\]

and

\[
\boxed{
\dim(\mathbb{R}^3)=3.
}
\]

\vspace{0.5cm}

More generally,

\[
\boxed{
\dim(\mathbb{R}^n)=n.
}
\]

Dimension tells us how many independent directions are needed
to describe the entire space.

\end{frame}


% =========================================================
% BASIS AND DIMENSION
% =========================================================

\begin{frame}{Basis, Span and Dimension}

These concepts are closely connected.
\vspace{-0.1cm}
\[
\boxed{
\begin{array}{c}
\text{Linearly independent vectors}\\[3pt]
\downarrow\\[3pt]
\text{Basis}\\[3pt]
\downarrow\\[3pt]
\text{Spans the vector space}\\[3pt]
\downarrow\\[3pt]
\text{Number of basis vectors}\\[3pt]
\downarrow\\[3pt]
\text{Dimension}
\end{array}
}
\]

A basis is therefore a \textbf{minimal set of generators} for
a vector space.

\end{frame}


% =========================================================
% VECTOR SPACE
% =========================================================

\begin{frame}{What Is a Vector Space?}

A \textbf{vector space} is a set of objects that is closed under:

\begin{itemize}
    \item vector addition
    \item scalar multiplication
\end{itemize}

and satisfies a collection of algebraic rules.

\vspace{0.4cm}

If $\vect{u}$ and $\vect{v}$ belong to a vector space $V$, and
$c$ is a scalar, then

\[
\vect{u}+\vect{v}\in V
\]

and

\[
c\vect{u}\in V.
\]

\vspace{0.4cm}

The scalars are usually elements of $\mathbb{R}$ or $\mathbb{C}$.

\end{frame}


% =========================================================
% VECTOR SPACE AXIOMS
% =========================================================

\begin{frame}{Vector Space Properties}

For vectors $\vect{u},\vect{v},\vect{w}\in V$ and scalar $c$,
the vector-space operations satisfy properties such as:

\begin{itemize}
    \item $\vect{u}+\vect{v}=\vect{v}+\vect{u}$
    \item $(\vect{u}+\vect{v})+\vect{w}
    =\vect{u}+(\vect{v}+\vect{w})$
    \item There exists a zero vector $\mathbf{0}$
    \item Every vector has an additive inverse
    \item $c(\vect{u}+\vect{v})
    =c\vect{u}+c\vect{v}$
    \item $(c+d)\vect{u}
    =c\vect{u}+d\vect{u}$
\end{itemize}

\vspace{0.2cm}

These properties provide the algebraic structure required
for vector-space operations.

\end{frame}


% =========================================================
% EXAMPLES OF VECTOR SPACES
% =========================================================

\begin{frame}{Examples of Vector Spaces}

Vector spaces are not limited to arrows in $\mathbb{R}^n$.

Examples include:

\begin{columns}[T,totalwidth=\textwidth]

\begin{column}{0.46\textwidth}

\begin{block}{$\mathbb{R}^n$}
\[
\begin{bmatrix}
x_1\\
\vdots\\
x_n
\end{bmatrix}
\]
\end{block}

\begin{block}{Polynomials}
\[
p(x)=a_0+a_1x+\cdots+a_nx^n
\]
\end{block}

\end{column}

\begin{column}{0.46\textwidth}

\begin{block}{Matrices}
Matrices of a fixed size form a vector space under
matrix addition and scalar multiplication.
\end{block}

\begin{block}{Functions}
Suitable sets of functions can also form vector spaces.
\end{block}

\end{column}

\end{columns}

\end{frame}


% =========================================================
% SCALARS
% =========================================================

\begin{frame}{Scalars and Fields}

A \textbf{scalar} is a quantity used to multiply vectors.

Common scalar sets include

\[
\mathbb{R}
\qquad\text{and}\qquad
\mathbb{C}.
\]

\vspace{0.4cm}

A \textbf{field} provides the arithmetic structure required
for scalar operations.

\vspace{0.4cm}

For example, the real numbers $\mathbb{R}$ form a field with

\begin{itemize}
    \item addition
    \item subtraction
    \item multiplication
    \item division by non-zero numbers
\end{itemize}

\vspace{0.2cm}

Vector spaces are therefore defined relative to a chosen
field of scalars.

\end{frame}


% =========================================================
% ALGEBRAIC STRUCTURES
% =========================================================

\begin{frame}{Algebraic Structures}

The lecture material introduces increasingly structured
algebraic systems.

\vspace{0.3cm}

\begin{center}

\begin{tikzpicture}[
    node distance=1.5cm,
    >=Latex
]

\node[
    draw,
    rounded corners,
    fill=LightBlue,
    minimum width=2.5cm
] (semigroup) {Semigroup};

\node[
    draw,
    rounded corners,
    fill=LightBlue,
    minimum width=2.5cm,
    right=of semigroup
] (monoid) {Monoid};

\node[
    draw,
    rounded corners,
    fill=LightLavender,
    minimum width=2.5cm,
    right=of monoid
] (group) {Group};

\node[
    draw,
    rounded corners,
    fill=LightLavender,
    minimum width=2.5cm,
    right=of group
] (abelian) {Abelian Group};

\draw[->, thick] (semigroup) -- (monoid);
\draw[->, thick] (monoid) -- (group);
\draw[->, thick] (group) -- (abelian);

\end{tikzpicture}

\end{center}

\vspace{0.4cm}

Each step adds additional algebraic properties.

\end{frame}


% =========================================================
% SEMIGROUP AND MONOID
% =========================================================

\begin{frame}{Semigroup and Monoid}

A \textbf{semigroup} is a set with an associative binary operation.

\[
(a*b)*c=a*(b*c).
\]

\vspace{0.4cm}

A \textbf{monoid} is a semigroup that additionally contains
an identity element $e$:

\[
\boxed{
e*a=a*e=a.
}
\]

\vspace{0.4cm}

These structures provide a foundation for understanding
more structured algebraic systems.

\end{frame}


% =========================================================
% GROUP
% =========================================================

\begin{frame}{Group and Abelian Group}

A \textbf{group} is a monoid in which every element has
an inverse.

For every $a$, there exists $a^{-1}$ such that

\[
aa^{-1}=a^{-1}a=e.
\]

\vspace{0.4cm}

An \textbf{Abelian group} additionally satisfies commutativity:

\[
\boxed{
a*b=b*a.
}
\]

\vspace{0.3cm}

Vector addition forms an Abelian group.

\end{frame}


% =========================================================
% VECTOR SPACE STRUCTURE
% =========================================================

\begin{frame}{From Algebraic Structure to Vector Spaces}

The vector-space structure combines two types of operations:

\vspace{0.4cm}

\begin{columns}[T,totalwidth=\textwidth]

\begin{column}{0.43\textwidth}

\begin{block}{Vector Addition}

\[
\vect{u}+\vect{v}
\]

Vectors form an Abelian group under addition.

\end{block}

\end{column}

\begin{column}{0.43\textwidth}

\begin{block}{Scalar Multiplication}

\[
c\vect{v}
\]

Scalars act on vectors while satisfying compatibility rules.

\end{block}

\end{column}

\end{columns}

\vspace{0.5cm}

Together, these operations produce the structure of a
\textbf{vector space}.

\end{frame}


% =========================================================
% VECTOR SPACES IN COMPUTING
% =========================================================

\begin{frame}{Vector Spaces in Computing}

The idea of a vector space allows many different objects
to be treated using the same mathematical framework.

\vspace{0.4cm}

\begin{itemize}
    \item Feature vectors in machine learning
    \item Coordinate vectors in computer graphics
    \item Matrix spaces in numerical computing
    \item Polynomial representations
    \item Function spaces in mathematical modelling
\end{itemize}

\begin{center}

\[
\boxed{
\text{Different objects}
\quad\longrightarrow\quad
\text{Common linear structure}
}
\]

\end{center}

\end{frame}